\documentclass[11pt]{article}

\usepackage[utf8]{inputenc}
\usepackage[T1]{fontenc}
\usepackage{geometry}
\usepackage{amsmath}
\usepackage{amssymb}
\usepackage{booktabs}
\usepackage{hyperref}
\usepackage{xurl}
\usepackage{xcolor}
\usepackage{listings}
\usepackage{enumitem}
\usepackage{parskip}
\usepackage{graphicx}
\graphicspath{{../images/}}

\geometry{margin=1in}

\lstset{
  language=Java,
  basicstyle=\ttfamily\small,
  keywordstyle=\color{blue},
  commentstyle=\color{gray},
  stringstyle=\color{teal},
  showstringspaces=false,
  columns=fullflexible,
  frame=single,
  framerule=0.2pt,
  breaklines=true
}

\setlist[itemize]{topsep=0.3em,itemsep=0.2em}
\setlist[enumerate]{topsep=0.3em,itemsep=0.2em}

% Simple per-section source line.
\newcommand{\notesources}[1]{\par\noindent{\small\color{gray}\textit{Sources:} \sloppy #1\par}}

% Unit-wise source strings for this subject (used by slides/notes).
% Path is relative to the lecture `latex/` folder (the compilation CWD).
% OOPS with Java (CSEG1044) — unit-wise sources
%
% These are short, student-facing reference strings intended to be used with:
%   - \slidesources{...}  (in beamer slides)
%   - \notesources{...}   (in lecture notes)
%
% Style inspiration: other subjects in My-Subjects/ use semicolon-separated sources
% and include an "accessed" date for web documentation.

\newcommand{\OOPSUnitISources}{Schildt, \textit{Java: A Beginner's Guide} (10e), McGraw-Hill (Java fundamentals + OOP basics).; Oracle Java Tutorials: Getting Started + Language Basics + Classes and Objects (accessed \today).; Java SE API docs (e.g., \texttt{String}, \texttt{Scanner}) (accessed \today).}

\newcommand{\OOPSUnitIISources}{Schildt, \textit{Java: The Complete Reference} (12e), McGraw-Hill (classes, inheritance, packages, interfaces).; Bloch, \textit{Effective Java} (3e), Addison-Wesley, 2018 (design choices; interfaces vs abstract classes).; Oracle Java Tutorials: Classes and Objects + Interfaces + Packages (accessed \today).; Java SE API docs (e.g., \texttt{Comparable}, \texttt{Runnable}) (accessed \today).}

\newcommand{\OOPSUnitIIISources}{Schildt, \textit{Java: The Complete Reference} (12e) (nested/inner classes, exceptions, threads, I/O).; Oracle Java Tutorials: Nested Classes + Exceptions + Concurrency + Basic I/O (accessed \today).; Java SE API docs (e.g., \texttt{Thread}, \texttt{AutoCloseable}) (accessed \today).}

\newcommand{\OOPSUnitIVSources}{Schildt, \textit{Java: The Complete Reference} (12e) (generics, lambdas, Swing, JDBC).; Bloch, \textit{Effective Java} (3e) (generics + lambdas).; Oracle Java Tutorials: Generics + Lambda Expressions + Swing + JDBC Basics (accessed \today).; Java SE API docs (e.g., \texttt{java.util.function}, \texttt{javax.swing}, \texttt{java.sql}) (accessed \today).}

\newcommand{\OOPSUnitVSources}{Schildt, \textit{Java: The Complete Reference} (12e) (Collections Framework + wrapper classes).; Oracle Java docs: Collections Framework overview (accessed \today).; Java SE API docs (\texttt{java.util} collections; wrapper classes in \texttt{java.lang}) (accessed \today).}

\newcommand{\OOPSUnitVISources}{Git SCM: \textit{Pro Git} (2e) (version control workflow), accessed \today.; GitHub Docs (repositories, issues, pull requests), accessed \today.; Oracle Java Tutorials: Swing + JDBC (accessed \today).; JUnit 5 User Guide (accessed \today).; Apache Maven Project: Getting Started (accessed \today).; Gradle User Manual: Getting Started (accessed \today).; UML class diagrams: UML 2.x overview/spec (accessed \today).}

\newcommand{\OOPSMiscWhyInterfacesSources}{Schildt, \textit{Java: The Complete Reference} (12e) (interfaces).; Bloch, \textit{Effective Java} (3e) (prefer interfaces where appropriate).; Oracle Java Tutorials: Interfaces (accessed \today).; Java SE API docs (e.g., \texttt{Comparable}, \texttt{Runnable}) (accessed \today).}



\title{Object Oriented Programming (CSEG1044)\\Unit 06 -- Lecture 04 Notes\\Database Design + JDBC Layer}
\author{Tofik Ali}
\date{\today}

\begin{document}
\maketitle
\tableofcontents

\section{Lecture Overview}
Your Swing UI is not complete until data is saved permanently.
This lecture converts your entity model into a database schema and creates the JDBC/DAO layer.
Deliverable-3 is a big milestone: once DAO CRUD works, UI integration becomes easy.
\notesources{\OOPSUnitVISources}

\section{Syllabus Alignment (Unit 06)}
\textbf{Unit 06:} database design + JDBC persistence

\section{Learning Outcomes}
After this lecture, you should be able to:
\begin{itemize}
  \item Convert an ER/entity model into tables with keys and constraints.
  \item Implement JDBC DAO methods using \texttt{PreparedStatement} and \texttt{ResultSet} mapping.
  \item Complete Deliverable-3: DB schema + CRUD for 1--2 core entities.
\end{itemize}

\section{Key Terms (Quick)}
\begin{itemize}
  \item \textbf{Primary key (PK)}: unique identifier for a row.
  \item \textbf{Foreign key (FK)}: reference to PK in another table.
  \item \textbf{Constraint}: rule enforced by DB (NOT NULL, UNIQUE, FK).
  \item \textbf{Index}: speeds up search (usually on PK/FK columns).
  \item \textbf{DAO}: class that performs CRUD using JDBC.
  \item \textbf{Mapping}: converting a DB row (\texttt{ResultSet}) into a Java object.
\end{itemize}

\section{Detailed Notes}
\subsection{ER/entity model $\rightarrow$ relational schema}
Start with the entities from Deliverable-2.
For each entity:
\begin{itemize}
  \item create a table,
  \item choose a primary key,
  \item choose column types (\texttt{INT}, \texttt{VARCHAR}, \texttt{TIMESTAMP}),
  \item add constraints (NOT NULL, UNIQUE),
  \item add foreign keys for relationships.
\end{itemize}

\paragraph{Example mapping.} Booking has a resource $\rightarrow$ \texttt{bookings.resource\_id} is a foreign key referencing \texttt{resources.id}.

\subsection{Keys and constraints (minimum)}
\begin{itemize}
  \item Use integer auto-increment PKs or meaningful unique keys (like roll number) depending on domain.
  \item Use NOT NULL for mandatory fields.
  \item Use UNIQUE where duplicates are not allowed (e.g., resource name).
  \item Use FK constraints to preserve relationships.
\end{itemize}

\subsection{JDBC layer (DAO pattern)}
Recommended structure:
\begin{itemize}
  \item \texttt{DbUtil}: provides DB connection.
  \item \texttt{ResourceDao}: CRUD for resources.
  \item \texttt{BookingDao}: CRUD for bookings.
\end{itemize}
Use \texttt{PreparedStatement} for all user-provided values.
\texttt{try-with-resources} ensures connections and statements are closed properly.
\notesources{\OOPSUnitVISources}

\section{Worked Example: schema + DAO skeleton}
\subsection*{Schema (example)}
\begin{lstlisting}[language=]
-- resources table
CREATE TABLE resources (
  id INT AUTO_INCREMENT PRIMARY KEY,
  name VARCHAR(100) NOT NULL UNIQUE,
  type VARCHAR(50) NOT NULL,
  capacity INT NOT NULL
);

-- bookings table
CREATE TABLE bookings (
  id INT AUTO_INCREMENT PRIMARY KEY,
  resource_id INT NOT NULL,
  user_name VARCHAR(100) NOT NULL,
  start_time TIMESTAMP NOT NULL,
  end_time TIMESTAMP NOT NULL,
  status VARCHAR(20) NOT NULL,
  CONSTRAINT fk_booking_resource FOREIGN KEY (resource_id)
    REFERENCES resources(id)
);
\end{lstlisting}

\subsection*{DbUtil (connection helper)}
\begin{lstlisting}
import java.sql.Connection;
import java.sql.DriverManager;
import java.sql.SQLException;

public class DbUtil {
    private static final String URL = "jdbc:mysql://localhost:3306/crm";
    private static final String USER = "root";
    private static final String PASS = "password";

    public static Connection getConnection() throws SQLException {
        return DriverManager.getConnection(URL, USER, PASS);
    }
}
\end{lstlisting}

\subsection*{ResourceDao (CRUD skeleton)}
\begin{lstlisting}
import java.sql.*;
import java.util.*;

public class ResourceDao {
    public void insert(Resource r) throws SQLException {
        String sql = "INSERT INTO resources(name, type, capacity) VALUES(?,?,?)";
        try (Connection con = DbUtil.getConnection();
             PreparedStatement ps = con.prepareStatement(sql)) {
            ps.setString(1, r.getName());
            ps.setString(2, r.getType());
            ps.setInt(3, r.getCapacity());
            ps.executeUpdate();
        }
    }

    public List<Resource> findAll() throws SQLException {
        String sql = "SELECT id, name, type, capacity FROM resources";
        List<Resource> out = new ArrayList<>();
        try (Connection con = DbUtil.getConnection();
             PreparedStatement ps = con.prepareStatement(sql);
             ResultSet rs = ps.executeQuery()) {
            while (rs.next()) {
                out.add(new Resource(
                    rs.getInt("id"),
                    rs.getString("name"),
                    rs.getString("type"),
                    rs.getInt("capacity")
                ));
            }
        }
        return out;
    }
}
\end{lstlisting}

\section{Deliverable-3 checklist}
\begin{itemize}
  \item ER/schema diagram or table list (resources, bookings, users)
  \item SQL schema script committed in repo (\texttt{schema.sql})
  \item DAO CRUD for 1--2 core entities (at least: insert + list)
  \item Simple test runner/main method that demonstrates DAO works
\end{itemize}

\section{In-Class Activity (evidence)}
\begin{itemize}
  \item Commit \texttt{schema.sql} and DAO code.
  \item Create PR for DB module.
  \item Add issues for remaining DAO methods (update/delete/search).
\end{itemize}

\section{Checkpoint Questions (with sample answers)}
\textbf{Q1.} Why use foreign keys?\par
\textbf{Sample answer:} Foreign keys enforce relationships between tables, preventing invalid references (e.g., a booking referencing a resource that does not exist).

\vspace{0.6em}
\textbf{Q2.} Why use \texttt{PreparedStatement}?\par
\textbf{Sample answer:} It supports parameterized queries, improves safety for user input, and avoids SQL injection risks compared to concatenating strings.

\end{document}
